\SourceLeaf{420}%
evidently have
\[
\begin{aligned}
\delta x&=\delta\varrho.\cos(\varphi+90^{0})\\
\delta y&=\delta\varrho.\sin(\varphi+90^{0})
\end{aligned}
\]
or
\[
\begin{aligned}
\delta x&=-\delta\varrho.\sin\varphi\\
\delta y&=+\delta\varrho.\cos\varphi.
\end{aligned}
\]
Now, when $\delta\varrho$ is infinitesimal, we have
\[
\begin{aligned}
\delta V&=p\delta x+q\delta y\\
&=(-p\sin\varphi+q\cos\varphi)\delta\varrho\\
&=\mp\delta\varrho\BareRoot{(pp+qq)}
\end{aligned}
\]
while on the curve itself $V$ $=0$. Thus the upper signs apply if $V$ passes from positive to negative when crossing the curve from left to right, and conversely. Combining this with what was said at the end of Article 2, it follows that the curve is always convex towards the side on which $V$ takes on the sign that
\[
Pqq-2Qpq+Rpp
\]
has.

If, for example, the curve is a circle and we put
\[
V=xx+yy-aa,
\]
then
\[
\begin{gathered}
p=2x,\qquad q=2y,\\
P=2,\qquad Q=0,\qquad R=2,\\
Pqq-2Qpq+Rpp=8yy+8xx=8aa,\\
(pp+qq)^{\frac32}=8a^3,\\
r=\pm a
\end{gathered}
\]
and the curve is convex towards the side on which
\[
xx+yy>aa,
\]
as required.
